% DONO: Kauê (+ números do Nathan)   ·   ESPAÇO: ~1,5 página — A SEÇÃO DE MAIOR PESO
\section{Avaliação}
% Opção B do enunciado: a verificação semântica funciona?
%
% Método:
% - N programas rotulados em avaliacao/casos/ (meta: 20 a 30), cada um com
%   gabarito "aceita" ou "recusa <categoria>".
% - Separar SEMENTES (13 casos gerados com IA, que já incluíam falhas
%   conhecidas) dos CASOS DA EQUIPE (escritos à mão, sem olhar o código).
%   Isso precisa estar explícito no texto.
% - Métricas: detectados, escapes, falsos positivos, aceitos corretamente;
%   tempo médio de verificação.
%
% Resultados: copiar a tabela gerada por `python avaliacao/rodar.py`.
% Resultado atual só das sementes (27/10), como referência:
%   detecção 6/7 (86%), falsos positivos 2/6 (33%), 1 escape, ~0,65 ms.
%
% Interpretação (o slide 20 exige):
% - De onde vêm os falsos positivos (c04: forma não normalizada; c05:
%   integralidade ignorada) e o escape (m02).
% - Antes x depois das melhorias do Nathan.
% - O que os números NÃO permitem concluir: os casos não vêm de lojas reais,
%   a amostra é pequena, os autores escreveram os casos e o verificador.
\begin{table}[h]
  \caption{Resultados da verificação semântica (preencher)}
  \begin{tabular}{lrrrr}
    \toprule
    categoria & casos & detectado & escape & falso pos. \\
    \midrule
    conflito\_desconto & & & & \\
    conflito\_frete & & & & \\
    regra\_morta & & & & \\
    controles (aceita) & & & & \\
    \bottomrule
  \end{tabular}
\end{table}
